% ============================================================
\chapter{How to use this book}\label{howto:chapter}
% ============================================================
% Template: replace every @@...@@ block; delete rows/sections for features switched off in PLAN.md.

\section*{The exam}
\begin{tabularx}{\linewidth}{@{}lX@{}}
\toprule
@@EXAM_ROWS@@
\bottomrule
\end{tabularx}

\section*{Structure}
@@STRUCTURE_PARAGRAPH@@
Each chapter opens with learning objectives and ends with take-home messages, a self-test and answers.
The overview map on the next pages shows how the chapters build on each other.

\section*{Box legend}
\begin{examcritical}
Facts and definitions you must be able to reproduce.
\end{examcritical}
\begin{tcolorbox}[studybox,colback=lecq!6,colframe=lecq,title={Lecturer question \#\hfill{\normalfont\footnotesize [Source page]}}]
\textbf{\itshape ``A question highlighted by the lecturer in the material.''}\par\smallskip
\textbf{Model answer.}\ A worked answer as you could write it in the exam.
\end{tcolorbox}
\begin{pitfall}
Typical confusions between similar concepts.
\end{pitfall}
\begin{keyidea}
The one idea a section hinges on.
\end{keyidea}
\begin{beyondslides}
Context that is \emph{not} in the course material (textbooks, general knowledge). Useful for
understanding, lower priority for memorising.
\end{beyondslides}

\section*{Source tags and how far to trust them}
Grey tags such as \src{@@EXAMPLE_TAG@@}{12} point to the \emph{PDF page} of the original material:

{\small\begin{tabularx}{\linewidth}{@{}l>{\raggedright\arraybackslash}X@{}}
\toprule
Tag & Source file\\
\midrule
@@TAG_ROWS@@
\bottomrule
\end{tabularx}}

A tag says \emph{where} a statement comes from, not that it is correct. Numbers and formulas were
matched against the text of the cited page; values that could not be confirmed (for example from
scanned pages) are printed like \unsure{this} and listed in the corrections appendix.
When a value matters for the exam, check it on the cited page.
